\documentclass[11pt, a4paper]{article}

% === LuaLaTeX 패키지 ===
\usepackage{fontspec}
\usepackage{kotex}
\usepackage{emoji}
\setemojifont{Noto Color Emoji}
\directlua{
  luaotfload.add_fallback("emojifb", {
    "NotoColorEmoji:mode=harf;"
  })
}
\setmainfont{Noto Sans CJK KR}[RawFeature={fallback=emojifb}]
\setsansfont{Noto Sans CJK KR}[RawFeature={fallback=emojifb}]
\setmonofont{Noto Sans Mono CJK KR}[Scale=0.85]

% === 레이아웃 ===
\usepackage[margin=2.3cm, top=2.5cm, bottom=2.5cm]{geometry}
\usepackage{setspace}\onehalfspacing

% === 패키지 ===
\usepackage{graphicx}
\usepackage{booktabs}
\usepackage{array}
\usepackage{longtable}
\usepackage{xcolor}
\usepackage{tcolorbox}
\usepackage{enumitem}
\usepackage{fancyhdr}
\usepackage{titlesec}
\usepackage{hyperref}
\usepackage{float}
\usepackage{fancyvrb}

% === 한글 캡션 ===
\renewcommand{\contentsname}{목차}
\renewcommand{\tablename}{표}
\renewcommand{\figurename}{그림}

% === 색상 ===
\definecolor{mainblue}{RGB}{21, 101, 192}
\definecolor{accentgreen}{RGB}{46, 125, 50}
\definecolor{warningorange}{RGB}{230, 81, 0}
\definecolor{lightbg}{RGB}{245, 245, 245}
\definecolor{tipbg}{RGB}{232, 245, 233}
\definecolor{warnbg}{RGB}{255, 243, 224}
\definecolor{corebg}{RGB}{227, 242, 253}

% === tcolorbox ===
\tcbuselibrary{skins, breakable}
\newtcolorbox{tipbox}[1][]{
  colback=tipbg, colframe=accentgreen, boxrule=1pt, arc=4pt,
  fonttitle=\bfseries\color{accentgreen}, title={#1}, breakable
}
\newtcolorbox{warnbox}[1][]{
  colback=warnbg, colframe=warningorange, boxrule=1pt, arc=4pt,
  fonttitle=\bfseries\color{warningorange}, title={#1}, breakable
}
\newtcolorbox{corebox}[1][]{
  colback=corebg, colframe=mainblue, boxrule=1pt, arc=4pt,
  fonttitle=\bfseries\color{mainblue}, title={#1}, breakable
}
\newtcolorbox{promptbox}{
  colback=lightbg, colframe=gray!60, boxrule=0.5pt, arc=2pt,
  fontupper=\ttfamily\small, breakable
}

% === 헤더/푸터 ===
\pagestyle{fancy}\fancyhf{}
\fancyhead[L]{\small\textcolor{mainblue}{사물인터넷과 빅데이터}}
\fancyhead[R]{\small\textcolor{gray}{2주차: 강의노트}}
\fancyfoot[C]{\thepage}
\renewcommand{\headrulewidth}{0.5pt}

% === 섹션 스타일 ===
\titleformat{\section}{\Large\bfseries\color{mainblue}}{\thesection.}{0.5em}{}
\titleformat{\subsection}{\large\bfseries\color{mainblue!80}}{\thesubsection}{0.5em}{}
\titleformat{\subsubsection}{\normalsize\bfseries\color{mainblue!60}}{\thesubsubsection}{0.5em}{}
\hypersetup{colorlinks=true, linkcolor=mainblue, urlcolor=mainblue!70}

% ============================================================
\begin{document}

% === 표지 ===
\begin{titlepage}
\centering
\vspace*{3cm}
{\Huge\bfseries\color{mainblue} 사물인터넷과 빅데이터\par}
\vspace{0.5cm}
{\LARGE 2주차 강의노트\par}
\vspace{1.5cm}
{\Huge\bfseries 바이브코딩 소개\par}
\vspace{2cm}
{\large 청주교육대학교 컴퓨터교육과\par}
\vspace{0.3cm}
{\large 2026학년도 1학기\par}
\vspace{3cm}
{\normalsize 문서 버전: v1.0.0.0\par}
{\normalsize 작성일: 2026-02-19\par}
\end{titlepage}

% === 목차 ===
\tableofcontents
\newpage

% ============================================================
\section{지난 주 돌아보기}
% ============================================================

\subsection{1주차 과제 점검}

지난 시간에 마이크로비트 LED에 자기 이름을 표시하는 과제를 했어요. 간단히 점검해 볼게요.

\begin{itemize}[leftmargin=1.5em]
\item MakeCode 에디터에 접속해서 블록 코딩을 해봤나요?
\item ``시작하면'' 블록 안에 ``문자열 표시'' 블록을 넣어서 이름을 표시했나요?
\item 시뮬레이터에서 LED가 스크롤되면서 이름이 나타나는 걸 확인했나요?
\end{itemize}

잘 됐든, 잘 안 됐든 괜찮아요. 오늘부터는 블록 코딩 대신 \textbf{완전히 새로운 방식}으로 프로그래밍을 해볼 거예요.

\subsection{오늘 배울 내용 미리보기}

오늘 수업은 크게 두 가지예요.

첫째, \textbf{바이브코딩}이라는 새로운 코딩 방식을 배워요. AI에게 ``이런 걸 만들어줘''라고 말하면 코드가 뚝딱 나오는 방법이에요.

둘째, \textbf{센서 탐색}을 해요. 마이크로비트에 내장된 센서와 추가로 연결할 수 있는 Grove 센서를 살펴보면서, ``우리 프로젝트에서 뭘 측정할 수 있을까?''를 생각해 볼 거예요.

% ============================================================
\section{바이브코딩이란?}
% ============================================================

\subsection{전통 코딩 vs 바이브코딩}

여러분이 알고 있는 코딩은 아마 이런 모습일 거예요.

\begin{corebox}[전통 코딩 방식]
프로그래밍 언어 공부 → 문법 암기 → 코드 작성 → 오류 수정 → 완성
\end{corebox}

변수, 함수, 반복문, 조건문... 이런 문법을 하나하나 배우고, 한 글자라도 틀리면 에러가 나죠. 세미콜론 하나 빠뜨려도 프로그램이 안 돌아가는 경험, 혹시 해보셨나요?

\begin{corebox}[바이브코딩 방식]
만들고 싶은 것을 말로 설명 → AI가 코드 생성 → 결과 확인 → 수정 요청 → 완성
\end{corebox}

바이브코딩(Vibe Coding)은 \textbf{코드를 직접 짜지 않고, AI에게 자연어로 설명하여 코드를 만드는 방법}이에요. ``바이브(vibe)''는 분위기, 느낌이라는 뜻인데, 내가 원하는 느낌을 AI에게 전달하면 AI가 코드로 바꿔준다는 거예요.

\subsection{바이브코딩의 핵심 아이디어}

바이브코딩에서 가장 중요한 건 \textbf{``무엇을 만들고 싶은지 명확하게 설명하는 능력''}이에요.

기존 코딩에서는 ``어떻게(How) 만드는가''가 중요했어요. \texttt{for} 반복문을 써야 하는지, \texttt{while}을 써야 하는지, 배열 인덱스는 0부터 시작하는지... 이런 기술적 지식이 필요했죠.

바이브코딩에서는 ``무엇을(What) 만들고 싶은가''가 중요해요. ``교실 온도를 그래프로 보여주는 웹페이지를 만들어줘''처럼, 원하는 결과를 잘 설명하면 돼요.

\subsection{바이브코딩의 장점과 한계}

\begin{tipbox}[장점]
\begin{itemize}[leftmargin=1.5em]
\item 프로그래밍 언어를 몰라도 프로그램을 만들 수 있어요
\item 아이디어를 빠르게 시제품(프로토타입)으로 만들 수 있어요
\item 코드 문법 대신 문제 해결 과정에 집중할 수 있어요
\item 결과물을 보면서 ``이건 좀 바꿔줘''라고 점진적으로 개선할 수 있어요
\end{itemize}
\end{tipbox}

\begin{warnbox}[한계]
\begin{itemize}[leftmargin=1.5em]
\item AI가 만든 코드가 항상 완벽하지는 않아요
\item 복잡한 프로그램일수록 원하는 대로 만들기 어려워져요
\item AI가 코드를 ``이해''하는 게 아니라 패턴을 기반으로 생성하는 거예요
\item 결과물을 검토하고 판단하는 건 결국 사람의 몫이에요
\end{itemize}
\end{warnbox}

\begin{corebox}[핵심 메시지]
``AI는 코드를 이해하지 않아요. 우리가 만든 설명을 바탕으로 코드를 생성할 뿐이에요. 그래서 결과를 꼭 확인해야 해요!''
\end{corebox}

\subsection{이 수업에서 바이브코딩을 사용하는 이유}

이 수업의 목표는 코딩 실력을 키우는 게 아니에요. \textbf{데이터를 수집하고, 분석하고, 해석하는 과정}을 경험하는 거예요.

바이브코딩을 사용하면 코드 문법에 시간을 쏟지 않고, 진짜 중요한 것들에 집중할 수 있어요.

\begin{itemize}[leftmargin=1.5em]
\item ``어떤 센서로 뭘 측정할까?'' (문제 정의)
\item ``수집한 데이터에서 어떤 패턴이 보일까?'' (데이터 분석)
\item ``이 패턴이 정말 원인--결과 관계일까?'' (비판적 사고)
\end{itemize}

코딩은 AI에게 맡기고, 우리는 \textbf{생각하는 일}에 집중하는 거예요.

% ============================================================
\section{Claude 사용법}
% ============================================================

\subsection{Claude 접속하기}

Claude는 Anthropic이라는 회사에서 만든 AI 도구예요. 웹 브라우저에서 바로 사용할 수 있어요.

\begin{enumerate}[leftmargin=1.5em]
\item 웹 브라우저에서 \texttt{claude.ai}에 접속해요
\item 학교에서 제공하는 계정으로 로그인해요
\item ``새 대화'' 버튼을 눌러 대화를 시작해요
\end{enumerate}

\subsection{Claude 화면 구성}

Claude 화면은 크게 세 영역으로 나뉘어요.

\begin{itemize}[leftmargin=1.5em]
\item \textbf{왼쪽 사이드바}: 이전 대화 목록이 있어요. 나중에 다시 찾아볼 수 있어요.
\item \textbf{가운데 대화 영역}: 여러분이 입력한 메시지와 Claude의 답변이 표시돼요.
\item \textbf{입력창}: 화면 하단에 있어요. 여기에 프롬프트를 입력해요.
\end{itemize}

\subsection{프롬프트 작성법}

프롬프트(Prompt)란 AI에게 보내는 요청 메시지예요. 프롬프트를 잘 쓰면 원하는 결과를 더 정확하게 얻을 수 있어요.

\subsubsection{원칙 1: 구체적으로 써요}

\begin{table}[H]\centering
\caption{모호한 vs 구체적인 프롬프트}
\begin{tabular}{p{4.5cm}p{8cm}}\toprule
\textbf{모호한 프롬프트} & \textbf{구체적인 프롬프트} \\\midrule
``웹페이지 만들어줘'' & ``내 이름, 취미, 좋아하는 과목을 소개하는 웹페이지를 만들어줘'' \\
``그래프 그려줘'' & ``온도 데이터를 시간순으로 꺾은선 그래프로 그려줘'' \\
\bottomrule
\end{tabular}
\end{table}

\subsubsection{원칙 2: 단계적으로 요청해요}

한 번에 모든 걸 요청하기보다, 단계별로 나눠서 요청하면 더 좋은 결과를 얻을 수 있어요.

\begin{promptbox}
{[1단계]} ``자기소개 홈페이지를 만들어줘. 이름은 홍길동이야.''\\
{[2단계]} ``배경색을 하늘색으로 바꿔줘.''\\
{[3단계]} ``좋아하는 과목 목록을 추가해줘.''
\end{promptbox}

\subsubsection{원칙 3: 예시를 포함해요}

원하는 결과의 예시를 보여주면 AI가 더 잘 이해해요.

\begin{promptbox}
``표를 만들어줘. 예를 들면 이런 식이야:\\
| 시간 | 온도 | 소음 |\\
| 9:00 | 24℃ | 45dB |''
\end{promptbox}

\subsection{Artifact 확인과 다운로드}

Claude가 코드를 만들면 \textbf{Artifact}라는 별도 패널에 결과가 표시돼요.

\begin{itemize}[leftmargin=1.5em]
\item \textbf{미리보기}: Artifact 패널에서 결과물을 바로 확인할 수 있어요
\item \textbf{코드 보기}: ``Code'' 탭을 누르면 실제 코드를 볼 수 있어요
\item \textbf{다운로드}: 우측 상단 다운로드 버튼으로 HTML 파일을 저장할 수 있어요
\end{itemize}

\subsection{대화목록 PDF 생성 방법}

과제 제출 시 대화 내역을 PDF로 제출해야 해요. Claude에게 다음 프롬프트를 입력하면 자동으로 생성돼요.

\begin{promptbox}
이 대화에서 사용한 프롬프트와 답변을 포함한 대화 전문과\\
프롬프트와 답변을 정리한 요약본을 PDF 문서로 만들어줘.\\
오늘 날짜, 수업주차 2주차, 작성자 이름 OOO, 학번 XXXXXX을 추가해줘.
\end{promptbox}

% ============================================================
\section{바이브코딩 체험 실습}
% ============================================================

\subsection{실습 과제: 자기소개 홈페이지 만들기}

이제 직접 바이브코딩을 해볼 거예요. Claude에게 자기소개 홈페이지를 만들어달라고 요청해 보세요.

\begin{promptbox}
나를 소개하는 웹페이지를 만들어줘.\\
- 이름: 홍길동\\
- 학교: 청주교육대학교 4학년\\
- 좋아하는 과목: 체육, 미술\\
- 취미: 독서, 영화 감상\\
- 배경색은 밝고 예쁘게 해줘
\end{promptbox}

\subsection{결과 확인 → 수정 요청}

Claude가 만든 결과물을 보고, 마음에 안 드는 부분이 있으면 수정을 요청해요.

\begin{promptbox}
``글씨가 좀 작아. 제목을 더 크게 해줘.''\\
``사진 넣을 자리를 추가해줘.''\\
``색깔을 좀 더 파스텔 톤으로 바꿔줘.''\\
``이모지를 추가해서 더 귀엽게 만들어줘.''
\end{promptbox}

이 과정이 바로 바이브코딩의 핵심이에요. \textbf{프롬프트 → 결과 확인 → 수정 요청}을 반복하면서 원하는 결과물에 점점 가까워지는 거예요.

\subsection{실습 후 생각해 볼 점}

\begin{itemize}[leftmargin=1.5em]
\item Claude가 내 의도를 잘 이해했나요?
\item 프롬프트를 어떻게 바꾸니까 결과가 달라졌나요?
\item Claude가 만든 코드에 혹시 이상한 부분은 없었나요?
\item 이 방법으로 더 복잡한 프로그램도 만들 수 있을까요?
\end{itemize}

\begin{corebox}[기억하세요]
``Claude가 만든 결과를 그대로 믿지 말고, 꼭 확인하는 습관을 들이세요!''
\end{corebox}

% ============================================================
\section{센서 탐색: 마이크로비트 내장 센서}
% ============================================================

\subsection{마이크로비트 V2에 들어있는 센서들}

지난 시간에 마이크로비트로 LED에 이름을 표시했죠? 사실 마이크로비트 안에는 \textbf{여러 종류의 센서}가 들어있어요. 추가 장치를 연결하지 않아도 바로 사용할 수 있는 센서들이에요.

\begin{table}[H]\centering
\caption{마이크로비트 V2 내장 센서}
\small
\begin{tabular}{llll}\toprule
\textbf{센서} & \textbf{측정 대상} & \textbf{값 범위} & \textbf{특징} \\\midrule
온도 센서 & 온도 (\textcelsius) & $-5$ \textasciitilde{} 50\textcelsius & CPU 근처, 실제보다 2\textasciitilde5\textcelsius{} 높을 수 있음 \\
마이크 & 소리 크기 & 0 \textasciitilde{} 255 & V2에서 새로 추가, 상대값 \\
조도 센서 & 밝기 & 0 \textasciitilde{} 255 & LED를 역이용하여 감지 \\
가속도계 & 움직임, 기울기 & $\pm$2g (3축) & 흔들림, 낙하 감지 \\
나침반 & 방향 & 0\textdegree{} \textasciitilde{} 360\textdegree & 지구 자기장 감지 \\
터치 로고 & 터치 여부 & 터치/비터치 & V2에서 새로 추가 \\
\bottomrule
\end{tabular}
\end{table}

\subsection{내장 센서의 특징과 한계}

\textbf{온도 센서의 특이한 점:} 마이크로비트의 온도 센서는 \textbf{프로세서(CPU) 안에 들어있는 온도 센서}예요. CPU가 작동하면서 발생하는 열 때문에 실제 공기 온도보다 2\textasciitilde5\textcelsius{} 높게 측정될 수 있어요.

\textbf{마이크의 특이한 점:} 마이크가 측정하는 값은 0\textasciitilde255 사이의 \textbf{상대적인 숫자}예요. ``소리 크기 = 150''이라고 해서 150데시벨이라는 뜻이 아니에요.

\textbf{조도 센서의 특이한 점:} 마이크로비트에는 별도 조도 센서가 없어요. 대신 \textbf{5\texttimes5 LED를 거꾸로 사용}해서 빛을 감지해요.

\subsection{내장 센서로 할 수 있는 것 vs 없는 것}

\begin{table}[H]\centering
\caption{내장 센서의 가능/불가능}
\begin{tabular}{p{6cm}p{6.5cm}}\toprule
\textbf{할 수 있는 것} & \textbf{할 수 없는 것} \\\midrule
실내 온도 변화 추적 & 정확한 기온 측정 \\
소음 수준 비교 (시끄러움/조용함) & 정확한 데시벨(dB) 측정 \\
밝음/어두움 구분 & 정확한 조도(lux) 측정 \\
움직임/기울기 감지 & 습도, 거리, 공기질 측정 \\
\bottomrule
\end{tabular}
\end{table}

% ============================================================
\section{센서 탐색: Grove 확장 센서}
% ============================================================

\subsection{Grove 시스템이란?}

Grove는 Seeed Studio라는 회사에서 만든 \textbf{플러그앤플레이 센서 시스템}이에요. 전용 케이블을 꽂기만 하면 끝이에요.

\begin{table}[H]\centering
\caption{전통 방식 vs Grove 방식}
\begin{tabular}{ll}\toprule
\textbf{전통 방식} & \textbf{Grove 방식} \\\midrule
브레드보드 + 점퍼선 + 저항 & 전용 케이블 한 개 \\
회로 지식 필요 & 회로 지식 불필요 \\
잘못 연결 시 고장 위험 & 잘못 끼울 수가 없는 구조 \\
연결에 10\textasciitilde30분 & 연결에 10초 \\
\bottomrule
\end{tabular}
\end{table}

\subsection{Grove Shield와 마이크로비트 연결}

\begin{enumerate}[leftmargin=1.5em]
\item 마이크로비트를 Grove Shield의 슬롯에 끼워요
\item Grove 센서를 Shield의 I2C 포트에 케이블로 연결해요
\item USB로 컴퓨터에 연결하면 끝!
\end{enumerate}

\subsection{I2C 통신이란?}

Grove 센서 대부분은 \textbf{I2C}(아이-투-씨)라는 통신 방식을 사용해요. 교실에 비유하면:

\begin{itemize}[leftmargin=1.5em]
\item \textbf{마이크로비트} = 선생님 (질문을 하는 사람)
\item \textbf{센서들} = 학생들 (대답하는 사람들)
\item \textbf{I2C 버스} = 인터콤 선 (모두 같은 선에 연결)
\item \textbf{I2C 주소} = 학생 번호 (각 센서의 고유 번호)
\end{itemize}

선생님이 ``3번 학생, 지금 온도가 몇 도야?''라고 물으면, 3번 학생만 대답해요. 이렇게 \textbf{하나의 선에 여러 센서를 연결}할 수 있어요.

\begin{table}[H]\centering
\caption{센서별 I2C 주소}
\begin{tabular}{ccl}\toprule
\textbf{센서} & \textbf{I2C 주소} & \textbf{역할} \\\midrule
SHT40 & 0x44 & 온도·습도 \\
TSL2561 & 0x29 & 조도 (정밀) \\
TCS34725 & 0x29 & 색상 \\
VL53L0X & 0x29 & 거리 (레이저) \\
BMP280 & 0x76 & 기압·고도 \\
BME680 & 0x76 & 환경 (온도·습도·기압·공기질) \\
LIS3DHTR & 0x18 & 가속도 (정밀) \\
\bottomrule
\end{tabular}
\end{table}

\begin{warnbox}[주소 충돌 주의]
같은 주소를 가진 센서들은 동시에 연결할 수 없어요. 예를 들어 TSL2561과 TCS34725는 둘 다 0x29이라 하나만 사용할 수 있어요.
\end{warnbox}

\subsection{주요 Grove I2C 센서 소개}

\textbf{SHT40} — 온도·습도 센서. 온도 $-40$\textasciitilde$125$\textcelsius($\pm$0.2\textcelsius), 습도 0\textasciitilde100\%($\pm$1.8\%). 마이크로비트 내장 온도보다 훨씬 정확해요.

\textbf{TSL2561} — 조도 센서. 0.1\textasciitilde40,000 lux. 사람 눈과 비슷한 감도, 정확한 lux 값을 제공해요.

\textbf{TCS34725} — 색상 센서. RGB 색상값 + 밝기를 측정해요.

\textbf{VL53L0X} — 거리 센서 (ToF). 30\textasciitilde2,000mm. 레이저를 쏘고 돌아오는 시간을 측정해요.

\textbf{BMP280} — 기압 센서. 300\textasciitilde1,100 hPa. 기압으로 고도 추정도 가능해요.

\textbf{BME680} — 환경 종합 센서. 온도, 습도, 기압, 공기질(VOC) 4가지를 한 번에 측정하는 만능 센서예요.

\textbf{LIS3DHTR} — 가속도 센서. $\pm$2/4/8/16g. 마이크로비트 내장보다 정밀하고 설정 범위가 넓어요.

\subsection{센서 조합의 다양성}

마이크로비트 내장 센서 4종(온도, 소음, 빛, 가속도)과 Grove 외부 센서를 조합하면, \textbf{848가지 이상의 센서 조합}이 가능해요.

\begin{table}[H]\centering
\caption{센서 조합 예시}
\begin{tabular}{lll}\toprule
\textbf{내장 센서} & \textbf{Grove 센서} & \textbf{측정 항목} \\\midrule
온도, 소음, 빛 & SHT40 (습도) & 온도·습도·소음·밝기 4가지 \\
소음 & BME680 (환경 종합) & 소음·온도·습도·기압·공기질 5가지 \\
가속도 & VL53L0X (거리) & 움직임·거리(사람 감지) 2가지 \\
\bottomrule
\end{tabular}
\end{table}

% ============================================================
\section{센서 시뮬레이터 소개}
% ============================================================

\subsection{센서 시뮬레이터 v9}

센서 시뮬레이터는 \textbf{실제 하드웨어 없이 센서 데이터를 체험}할 수 있는 웹 도구예요. 마이크로비트와 Grove Shield를 화면에서 보여주고, 최대 6개의 센서를 배치하면서 가상 데이터를 확인할 수 있어요.

\subsection{센서 검색 도구}

``교실 온도와 습도를 측정하고 싶어요''라고 입력하면, 사용 가능한 내장 센서와 외부 센서를 추천해 줘요.

\subsection{I2C 연결도}

센서별 I2C 주소와 연결 구조를 시각적으로 보여주는 도구예요. 같은 주소를 가진 센서끼리는 연결할 수 없다는 것도 이 도구로 확인할 수 있어요.

% ============================================================
\section{프로젝트 브레인스토밍 안내}
% ============================================================

\subsection{다음 주 예고: 프로젝트 기획}

다음 주(3주차)에는 여러분만의 \textbf{프로젝트 주제}를 정할 거예요.

\begin{table}[H]\centering
\caption{프로젝트 예시}
\begin{tabular}{lll}\toprule
\textbf{주제} & \textbf{센서} & \textbf{핵심 질문} \\\midrule
교실 환경과 집중도 & 온도, 소음, 빛, 습도 & 교실 환경이 집중도에 영향을 줄까? \\
복도 소음 분석 & 소음 & 쉬는 시간 vs 수업 시간, 소음 차이는? \\
급식실 대기 시간 & 거리 (사람 감지) & 줄이 가장 긴 시간대는? \\
운동장 활동량 & 가속도 & 체육 시간에 얼마나 움직일까? \\
도서관 이용 패턴 & 빛, 소음 & 가장 조용한 시간대는? \\
\bottomrule
\end{tabular}
\end{table}

\subsection{다음 주까지 생각해 올 것}

\begin{itemize}[leftmargin=1.5em]
\item 내가 실습학교에서 측정하고 싶은 것은?
\item 어떤 센서가 필요할까?
\item 데이터를 모으면 어떤 질문에 답할 수 있을까?
\end{itemize}

% ============================================================
\section{오늘 수업 정리}
% ============================================================

\subsection{핵심 정리}

\textbf{바이브코딩}: 코드를 직접 짜지 않고, AI에게 자연어로 설명하여 코드를 만드는 방법이에요. 프롬프트를 잘 쓰는 게 핵심이에요.

\textbf{Claude 사용법}: 구체적으로, 단계적으로, 예시를 포함해서 프롬프트를 작성하면 좋은 결과를 얻을 수 있어요.

\textbf{마이크로비트 내장 센서}: 온도, 소음, 빛, 가속도, 나침반, 터치 6가지 센서가 내장되어 있어요.

\textbf{Grove 확장 센서}: I2C 방식으로 플러그앤플레이 연결이 가능한 외부 센서예요. SHT40(온습도), TSL2561(조도), TCS34725(색상), VL53L0X(거리), BMP280(기압), BME680(환경), LIS3DHTR(가속도) 등을 추가할 수 있어요.

\subsection{과제 안내}

\begin{table}[H]\centering
\caption{과제 안내}
\begin{tabular}{ll}\toprule
\textbf{항목} & \textbf{내용} \\\midrule
과제명 & 바이브코딩 체험 + 센서 탐색 \\
내용 & Claude로 자기소개 홈페이지 만들기 \\
제출물 & 결과물(HTML) + 대화목록 전문(PDF) + 대화목록 요약본(PDF) \\
제출 기한 & 일요일 자정 (LMS) \\
배점 & 5\% \\
\bottomrule
\end{tabular}
\end{table}

% ============================================================
\section*{참고자료}
\addcontentsline{toc}{section}{참고자료}

\begin{table}[H]\centering
\begin{tabular}{clll}\toprule
\textbf{\#} & \textbf{자료명} & \textbf{유형} & \textbf{비고} \\\midrule
1 & 강의계획서 v1.5.3.0 & 문서 & 2주차 상세 계획 \\
2 & 2주차 강의개요 v1.0.0.0 & 문서 & 본 강의노트의 기반 \\
3 & 센서 시뮬레이터 v9.1.0.0 & HTML & 참고 도구 \\
4 & 센서 검색 도구 v1.0.2.0 & HTML & 참고 도구 \\
5 & I2C 연결도 v1.0.2.0 & HTML & 참고 도구 \\
\bottomrule
\end{tabular}
\end{table}

% ============================================================
\section*{생성/수정 이력}
\addcontentsline{toc}{section}{생성/수정 이력}

\begin{table}[H]\centering
\begin{tabular}{llllll}\toprule
\textbf{버전} & \textbf{날짜} & \textbf{시간} & \textbf{수준} & \textbf{변경 내용} & \textbf{작성자} \\\midrule
v1.0.0.0 & 2026-02-19 & 21:30 & 최초 & 강의개요 기반 2주차 강의노트 초안 & Claude \\
\bottomrule
\end{tabular}
\end{table}

\end{document}
